% GENERATED FILE. Edit canonical lessons, then run npm run generate:books.
% canonical-source-hash: fnv1a64:a2c086d96c8a83ec
% canonical-lessons: BN-C352-boishomyo, BN-C352-morjada, BN-C352-sohanubhuti, BN-C352-bonchito, BN-C352-ontorbhukto-kora

\chapter{Fair for Everyone}
\label{ch:bn-fair-for-everyone}
\hlchaptermodality{352}

\begin{chapteropening}
\textbf{By the end of this chapter:} \emph{I can talk about unfairness and who is left out.}

Complete the last lesson of chapter 352: I can talk about unfairness and who is left out.
\end{chapteropening}

\section[\texorpdfstring{\bn{বৈষম্য} (boishômyô)}{বৈষম্য (boishômyô)}]{\bn{বৈষম্য} (boishômyô) — discrimination; inequality}

\label{lesson:BN-C352-boishomyo}



\begin{quote}
Before the new one: say the Bengali for esteem, then the Bengali for to treat with regard.
\end{quote}

\subsection*{You'll want to know: \bn{বৈষম্য}}
\textbf{\bn{বৈষম্য}} — \emph{boishômyô} — \textquotedblleft{}discrimination; inequality\textquotedblright{}.

\begin{quote}
\bn{মেয়েদের} \bn{সঙ্গে} \bn{বৈষম্য} \bn{করা} \bn{অন্যায়।}
\end{quote}

\begin{quote}
\emph{meyeder sôngge boishômyô kôrā ônyāy.}
\end{quote}

\begin{quote}
Discriminating against girls is wrong.
\end{quote}

\begin{cousinweb}
\textbf{\bn{বৈষম্য}} (\emph{boishômyô}) is Sanskrit \emph{vaiṣamya}, unevenness (Monier-Williams), from \emph{viṣama} 'uneven'. Bengali uses it for unequal treatment; it is done with -\bn{এর} \bn{সঙ্গে}.
\end{cousinweb}

\begin{grammarlens}[title={What you've built}]
One more word for giving reasons and opinions.
\end{grammarlens}

\subsection*{Your turn}
\begin{itemize}
\raggedright
  \item \emph{Say it:} \emph{boishômyô}
  \item \emph{Say it:} \emph{boishômyô}, once more
  \item \emph{Say it:} \emph{meyeder sôngge boishômyô kôrā ônyāy.} — Discriminating against girls is wrong.
  \item \emph{Recall:} say the Bengali for esteem, then the Bengali for to treat with regard, then say \emph{boishômyô} again
\end{itemize}

\begin{tcolorbox}[breakable,colback=teal!4,colframe=teal!35!black,title={Before you move on}]
What does \bn{বৈষম্য} mean? (\textquotedblleft{}Discrimination; inequality\textquotedblright{}.) Say it once more.
\end{tcolorbox}
\section[\texorpdfstring{\bn{মর্যাদা} (môrjādā)}{মর্যাদা (môrjādā)}]{\bn{মর্যাদা} (môrjādā) — dignity; standing}

\label{lesson:BN-C352-morjada}



\begin{quote}
Before the new one: say the Bengali for to treat with regard, then the Bengali for discrimination.
\end{quote}

\subsection*{You'll want to know: \bn{মর্যাদা}}
\textbf{\bn{মর্যাদা}} — \emph{môrjādā} — \textquotedblleft{}dignity; standing\textquotedblright{}.

\begin{quote}
\bn{সব} \bn{মানুষের} \bn{মর্যাদা} \bn{সমান।}
\end{quote}

\begin{quote}
\emph{sôb mānusher môrjādā sômān.}
\end{quote}

\begin{quote}
Every person has equal dignity.
\end{quote}

\begin{cousinweb}
\textbf{\bn{মর্যাদা}} (\emph{môrjādā}) is Sanskrit \emph{maryādā}, a boundary and so the proper limits of conduct (Monier-Williams). Bengali uses it for the dignity or standing those limits protect.
\end{cousinweb}

\begin{grammarlens}[title={What you've built}]
One more word for giving reasons and opinions.
\end{grammarlens}

\subsection*{Your turn}
\begin{itemize}
\raggedright
  \item \emph{Say it:} \emph{môrjādā}
  \item \emph{Say it:} \emph{môrjādā}, once more
  \item \emph{Say it:} \emph{sôb mānusher môrjādā sômān.} — Every person has equal dignity.
  \item \emph{Recall:} say the Bengali for to treat with regard, then the Bengali for discrimination, then say \emph{môrjādā} again
\end{itemize}

\begin{tcolorbox}[breakable,colback=teal!4,colframe=teal!35!black,title={Before you move on}]
What does \bn{মর্যাদা} mean? (\textquotedblleft{}Dignity; standing\textquotedblright{}.) Say it once more.
\end{tcolorbox}
\section[\texorpdfstring{\bn{সহানুভূতি} (sôhānubhūti)}{সহানুভূতি (sôhānubhūti)}]{\bn{সহানুভূতি} (sôhānubhūti) — sympathy}

\label{lesson:BN-C352-sohanubhuti}



\begin{quote}
Before the new one: say the Bengali for discrimination, then the Bengali for dignity.
\end{quote}

\subsection*{You'll want to know: \bn{সহানুভূতি}}
\textbf{\bn{সহানুভূতি}} — \emph{sôhānubhūti} — \textquotedblleft{}sympathy\textquotedblright{}.

\begin{quote}
\bn{সবাই} \bn{মীরাকে} \bn{সহানুভূতি} \bn{জানাল।}
\end{quote}

\begin{quote}
\emph{sôbāi mīrāke sôhānubhūti jānālô.}
\end{quote}

\begin{quote}
Everyone offered Mira their sympathy.
\end{quote}

\begin{cousinweb}
\textbf{\bn{সহানুভূতি}} (\emph{sôhānubhūti}) is a later compound of Sanskrit \emph{saha}, with, and \emph{anubhūti}, perception or feeling (Monier-Williams): feeling along with someone.
\end{cousinweb}

\begin{grammarlens}[title={What you've built}]
One more word for giving reasons and opinions.
\end{grammarlens}

\subsection*{Your turn}
\begin{itemize}
\raggedright
  \item \emph{Say it:} \emph{sôhānubhūti}
  \item \emph{Say it:} \emph{sôhānubhūti}, once more
  \item \emph{Say it:} \emph{sôbāi mīrāke sôhānubhūti jānālô.} — Everyone offered Mira their sympathy.
  \item \emph{Recall:} say the Bengali for discrimination, then the Bengali for dignity, then say \emph{sôhānubhūti} again
\end{itemize}

\begin{tcolorbox}[breakable,colback=teal!4,colframe=teal!35!black,title={Before you move on}]
What does \bn{সহানুভূতি} mean? (\textquotedblleft{}Sympathy\textquotedblright{}.) Say it once more.
\end{tcolorbox}
\section[\texorpdfstring{\bn{বঞ্চিত} (bônchitô)}{বঞ্চিত (bônchitô)}]{\bn{বঞ্চিত} (bônchitô) — deprived (of)}

\label{lesson:BN-C352-bonchito}



\begin{quote}
Before the new one: say the Bengali for dignity, then the Bengali for sympathy.
\end{quote}

\subsection*{You'll want to know: \bn{বঞ্চিত}}
\textbf{\bn{বঞ্চিত}} — \emph{bônchitô} — \textquotedblleft{}deprived (of)\textquotedblright{}.

\begin{quote}
\bn{গ্রামের} \bn{মানুষ} \bn{ভালো} \bn{হাসপাতাল} \bn{থেকে} \bn{বঞ্চিত।}
\end{quote}

\begin{quote}
\emph{grāmer mānush bhālo hāspātāl theke bônchitô.}
\end{quote}

\begin{quote}
Villagers are deprived of a good hospital.
\end{quote}

\begin{cousinweb}
\textbf{\bn{বঞ্চিত}} (\emph{bônchitô}) is Sanskrit \emph{vañcita}, deceived or cheated (Monier-Williams). The Bengali sense, deprived of something, is not one the dictionary gives; what is missing takes \bn{থেকে}.
\end{cousinweb}

\begin{grammarlens}[title={What you've built}]
One more word for giving reasons and opinions.
\end{grammarlens}

\subsection*{Your turn}
\begin{itemize}
\raggedright
  \item \emph{Say it:} \emph{bônchitô}
  \item \emph{Say it:} \emph{bônchitô}, once more
  \item \emph{Say it:} \emph{grāmer mānush bhālo hāspātāl theke bônchitô.} — Villagers are deprived of a good hospital.
  \item \emph{Recall:} say the Bengali for dignity, then the Bengali for sympathy, then say \emph{bônchitô} again
\end{itemize}

\begin{tcolorbox}[breakable,colback=teal!4,colframe=teal!35!black,title={Before you move on}]
What does \bn{বঞ্চিত} mean? (\textquotedblleft{}Deprived (of)\textquotedblright{}.) Say it once more.
\end{tcolorbox}
\section[\texorpdfstring{\bn{অন্তর্ভুক্ত} \bn{করা} (ôntôrbhuktô kôrā)}{অন্তর্ভুক্ত করা (ôntôrbhuktô kôrā)}]{\bn{অন্তর্ভুক্ত} \bn{করা} (ôntôrbhuktô kôrā) — to include}

\label{lesson:BN-C352-ontorbhukto-kora}



\begin{quote}
Before the new one: say the Bengali for sympathy, then the Bengali for deprived.
\end{quote}

\subsection*{You'll want to know: \bn{অন্তর্ভুক্ত} \bn{করা}}
\textbf{\bn{অন্তর্ভুক্ত} \bn{করা}} — \emph{ôntôrbhuktô kôrā} — \textquotedblleft{}to include\textquotedblright{}.

\begin{quote}
\bn{আমাদের} \bn{দলে} \bn{মীরাকেও} \bn{অন্তর্ভুক্ত} \bn{করো।}
\end{quote}

\begin{quote}
\emph{āmāder dôle mīrākeo ôntôrbhuktô kôro.}
\end{quote}

\begin{quote}
Include Mira in our team too.
\end{quote}

Past (he/she …): \bn{তিনি} \bn{আমাকে} \bn{দলে} \bn{অন্তর্ভুক্ত} \bn{করলেন}, \emph{tini āmāke dôle ôntôrbhuktô kôrlen}.

\begin{cousinweb}
\textbf{\bn{অন্তর্ভুক্ত} \bn{করা}} (\emph{ôntôrbhuktô kôrā}) is a learned compound of Sanskrit \emph{antar}, within, with \emph{bhukta}, eaten or enjoyed (Monier-Williams), so 'taken in', and \bn{করা}: to take someone or something inside a group. The compound itself is not in the dictionary.
\end{cousinweb}

\begin{grammarlens}[title={What you've built}]
One more word for giving reasons and opinions.
\end{grammarlens}

\subsection*{Your turn}
\begin{itemize}
\raggedright
  \item \emph{Say it:} \emph{ôntôrbhuktô kôrā}
  \item \emph{Say it:} \emph{ôntôrbhuktô kôrā}, once more
  \item \emph{Say it:} \emph{āmāder dôle mīrākeo ôntôrbhuktô kôro.} — Include Mira in our team too.
  \item \emph{Recall:} say the Bengali for sympathy, then the Bengali for deprived, then say \emph{ôntôrbhuktô kôrā} again
\end{itemize}

\begin{tcolorbox}[breakable,colback=teal!4,colframe=teal!35!black,title={Before you move on}]
What does \bn{অন্তর্ভুক্ত} \bn{করা} mean? (\textquotedblleft{}To include\textquotedblright{}.) Say it once more.
\end{tcolorbox}
